% Internal research-in-progress draft; authors and affiliation are unconfirmed.
% Reference: D:/hoctap/hybird control system.pdf, supplied by the user.
% Elsevier CAS-based layout adapted to the measured production typography.
% No publisher identity, article DOI, acceptance dates or publication status
% from the sample is reproduced. This is not a journal-issued production proof.
\documentclass[fleqn]{cas-dc}
\usepackage[numbers,sort&compress]{natbib}
\usepackage{fontspec}
\setmainfont{Charis SIL}[
  Path=open_fonts/,
  UprightFont=CharisSIL-Regular.ttf,
  BoldFont=CharisSIL-Bold.ttf,
  ItalicFont=CharisSIL-Italic.ttf,
  BoldItalicFont=CharisSIL-BoldItalic.ttf
]
\geometry{paperwidth=210mm,paperheight=280mm,
  left=13.26mm,right=13.2mm,top=19.5mm,bottom=20mm,
  headsep=10pt,footskip=12pt,columnsep=18pt}
\hypersetup{colorlinks=true,citecolor=blue!55!black,
  urlcolor=blue!55!black,linkcolor=blue!55!black}
\urlstyle{same}
\makeatletter
\renewcommand\normalsize{\@setfontsize\normalsize{8}{10.5}}
\renewcommand\small{\@setfontsize\small{7.2}{9}}
\renewcommand\footnotesize{\@setfontsize\footnotesize{6.4}{8.5}}
\renewcommand\scriptsize{\@setfontsize\scriptsize{6}{8}}
\normalsize
\renewcommand\sectionfont{\rmfamily\fontsize{8}{10.5}\bfseries\selectfont}
\renewcommand\ssectionfont{\rmfamily\fontsize{8}{10.5}\itshape\selectfont}
\renewcommand\printorcid{}
\AtBeginDocument{\renewcommand\bibfont{\fontsize{7.2}{9.5}\selectfont}}
\setlength{\parindent}{1em}
\setlength{\parskip}{0pt}
\newcommand{\DTAbstractFontAudit}{%
  \typeout{DTCHECK abstract size=\f@size pt; baseline=\the\baselineskip;
    family=\f@family; series=\f@series; shape=\f@shape}%
}
\def\ps@draftnotice{%
  \def\@oddhead{\hfil{\footnotesize\itshape Research-in-progress manuscript}\hfil}%
  \let\@evenhead\@oddhead
  \def\@oddfoot{\hfil{\footnotesize\thepage}\hfil}%
  \let\@evenfoot\@oddfoot
}
\makeatother
% Adapt the public CAS front matter without editing the vendor class files.
% Fractions reproduce the supplied sample's narrow article-info panel.
\ExplSyntaxOn
\RenewDocumentEnvironment { Abstract } { o }
{
  \group_begin:
  \parindent0pt
  \box_if_empty:NTF \g_stm_key_box
    { \leftskip=.316\textwidth }
    {
      \dim_set:Nn \l_tmpa_dim { \box_ht:N\g_stm_key_box + \box_dp:N\g_stm_key_box }
      \leftskip=.316\textwidth
      \hspace*{-.316\textwidth}
      \noindent\hbox_to_wd:nn { 0pt } { \box\g_stm_key_box\hss }
      \skip_vertical:n { -\l_tmpa_dim }
    }
  \noindent{\fontsize{7}{9}\selectfont A\,B\,S\,T\,R\,A\,C\,T}\par
  \skip_vertical:n { -4pt }
  \noindent\rule{.684\textwidth}{.2pt}\par
  \fontsize{7.2}{9}\selectfont
  \DTAbstractFontAudit
  \ignorespaces
}
{ \par\group_end: }
\RenewDocumentEnvironment { keywords } { O{Keywords} }
{
  \tex_global:D\tex_setbox:D\g_stm_key_box=\vtop\bgroup
  \hsize=.256\textwidth
  \cs_set:Npn\sep{\par}
  \parindent0pt
  {\fontsize{7}{9}\selectfont A\,R\,T\,I\,C\,L\,E\ \ I\,N\,F\,O}\par
  \skip_vertical:n{-3pt}
  \rule{.256\textwidth}{.2pt}\par\footnotesize
  \noindent\textit{#1}:\par
}
{ \egroup }
\ExplSyntaxOff
\begin{document}
\shorttitle{Network-Aware Predictive Digital Twin}
\shortauthors{}
\title[mode=title,size={\fontsize{13.5}{17.2}\selectfont},after=15pt]
{Network-Aware Predictive Digital Twin Framework for Resilient and Self-Reconfigurable Cyber-Physical Production Systems}
% Blank, reserved fields: user has not yet supplied an author list or affiliation.
\author{\mbox{}}
\begin{abstract}
Sensor faults and delayed state information can cause a production controller and its digital representation to disagree about material flow and equipment availability. This study proposes a network-aware predictive digital twin framework for software-level reconfiguration of a conveyor-based sorting and pick-and-place cell. A MuJoCo model of two conveyors, a pneumatic transfer mechanism, a Cartesian manipulator, and a vacuum gripper provides the initial simulation basis. The proposed architecture separates timestamped sensor observations from simulator ground truth and combines hybrid state estimation, residual-based anomaly detection, and scenario rollouts constrained by a computation deadline. Communication delay, jitter, packet loss, and data freshness are treated as inputs to state confidence and supervisory decisions. Candidate responses include adjusting conveyor timing, restricting material admission, using the remaining valid observations, and holding operation when uncertainty exceeds a task-dependent limit. Reconfiguration is restricted to policies supported by the installed actuators and local control interlocks. The planned evaluation compares sequential control, fault-aware supervision, and network-aware predictive supervision under matched fault schedules and operating conditions. Assessment criteria include state-estimation error, fault-detection delay, false-alarm frequency, task completion, recovery time, and prediction-deadline compliance. The current prototype supplies the physical simulation environment; the estimator, network-aware supervisor, and their quantitative validation remain under development. The research objective is to determine when predictive supervision can preserve production continuity under sensing and communication disturbances without changing the mechanical configuration.
\end{abstract}
\begin{keywords}
Cyber-physical production systems \sep
Digital twin \sep
Fault diagnosis \sep
Network-aware control \sep
Software reconfiguration
\end{keywords}
\maketitle
\pagestyle{draftnotice}
\thispagestyle{draftnotice}
\normalsize
\end{document}

